% Backup created 2026-09-09. Removed from the main text when Section 4.3 (general two-source learning theory)
% was replaced by the fitting-and-coefficient-selection subsection. Not input by main.tex.

% ===== Former Section 4.3 =====
\subsection{General two-source learning theory}
\label{sec:cate-general-learning}

Fix a method $j\in\{\mathrm{DRF},\mathrm{RF}\}$, coefficients
$(\lambda,\omega)$, and honestly fitted nuisances.  Let $n_f$ and $N_f$
denote the independent RCT and OBS fitting-sample sizes.  The method-specific
losses in Equations~\eqref{eq:dr-fusion-loss} and
\eqref{eq:rf-fusion-loss} can be written as an RCT contribution
$\ell_{R,j}(\zeta)$ and an OBS contribution $\ell_{O,j}(\zeta)$.

\begin{definition}[General fusion empirical learner]
For $\zeta\in\mathcal Z$, define
\begin{align*}
L_j(\zeta)
&\triangleq P_R\ell_{R,j}(\zeta)+P_O\ell_{O,j}(\zeta),\\
\widehat L_j(\zeta)
&\triangleq
\mathbb P_{R,n_f}\ell_{R,j}(\zeta)
+\mathbb P_{O,N_f}\ell_{O,j}(\zeta).
\end{align*}
Let
\begin{align*}
\zeta_j^\circ
&\in\operatorname*{arg\,min}_{\zeta\in\mathcal Z}\mathcal R(\zeta)
\end{align*}
be a prespecified, deterministically selected best-in-class function.  Let
$\operatorname{Pen}:\mathcal Z\to[0,\infty]$ be the full penalty functional
and define
\begin{align*}
\widehat J_j(\zeta)
&\triangleq\widehat L_j(\zeta)+\operatorname{Pen}(\zeta).
\end{align*}
Assume $\widehat J_j$ is bounded below and has a measurable global
$\epsilon_{\mathrm{opt}}$-approximate minimizer $\widehat\zeta_j$, meaning
\begin{align*}
\widehat J_j(\widehat\zeta_j)
&\leq\inf_{\zeta\in\mathcal Z}\widehat J_j(\zeta)
+\epsilon_{\mathrm{opt}}.
\end{align*}
Finally, define the centered comparator process
\begin{align*}
\mathbb Z_j(\zeta)
&\triangleq
\{\widehat L_j-L_j\}(\zeta)
-\{\widehat L_j-L_j\}(\zeta_j^\circ).
\end{align*}
\end{definition}

The next theorem is deterministic once the two tuning samples and fitted
nuisances are fixed.  It is the master reduction from fusion-loss fluctuation
to CATE excess risk.
\begin{theorem}[Exact general fusion ERM reduction]
\label{thm:cate-general-erm}
Assume RCT consistency and conditional randomization, a known propensity
$e$ satisfying strict overlap, and honest nuisance fitting.  Suppose either
$P_R^X=P_O^X$ and $r_0=1$, or $P_R^X\ll P_O^X$ and the exact ratio
$r_0=dP_R^X/dP_O^X$ is used.  Suppose the displayed losses are finite and
$\widehat J_j$ is bounded below and has a measurable global
$\epsilon_{\mathrm{opt}}$-approximate minimizer $\widehat\zeta_j$, meaning
$\widehat J_j(\widehat\zeta_j)\leq
\inf_{\zeta\in\mathcal Z}\widehat J_j(\zeta)+\epsilon_{\mathrm{opt}}$. Then
\begin{align}
\mathcal R(\widehat\zeta_j)-\mathcal R(\zeta_j^\circ)
&\leq-\mathbb Z_j(\widehat\zeta_j)
+\operatorname{Pen}(\zeta_j^\circ)
-\operatorname{Pen}(\widehat\zeta_j)
+\epsilon_{\mathrm{opt}}.
\label{eq:cate-general-erm-exact}
\end{align}
Consequently, the right-hand side is at most
$\sup_{\zeta\in\mathcal Z}\{-\mathbb Z_j(\zeta)\}
+\operatorname{Pen}(\zeta_j^\circ)+\epsilon_{\mathrm{opt}}$.
\end{theorem}

To obtain a sample-size-dependent result, localize by excess risk rather than
by an unnamed norm.  For $u\geq0$, define
\begin{align*}
\mathcal Z_j(u)
&\triangleq\{\zeta\in\mathcal Z:
\mathcal R(\zeta)-\mathcal R(\zeta_j^\circ)\leq u\}.
\end{align*}
Let $\mathbb Z_{R,j}$ and $\mathbb Z_{O,j}$ be the sourcewise centered
increments whose sum is $\mathbb Z_j$.  Suppose measurable functions
$\psi_{R,j},\psi_{O,j}$ and $\operatorname{Tail}_{R,j},
\operatorname{Tail}_{O,j}$ satisfy, sourcewise and uniformly over the
localized class, the high-probability upper bounds stated in the theorem
below. Define
\begin{align*}
\Phi_{j,n_f,N_f}(u;\delta)
&\triangleq
\psi_{R,j}(u;n_f)+\psi_{O,j}(u;N_f)
+\operatorname{Tail}_{R,j}(u,\delta)
+\operatorname{Tail}_{O,j}(u,\delta),\\
u_{j,n_f,N_f}^\star(\delta)
&\triangleq
\inf\{u>0:\Phi_{j,n_f,N_f}(u;\delta)\leq c_0u\},
\qquad 0<c_0<1.
\end{align*}
For every fixed $\delta$, assume that the full function
$u\mapsto\Phi_{j,n_f,N_f}(u;\delta)$ is nondecreasing and that
$u\mapsto\Phi_{j,n_f,N_f}(u;\delta)/\sqrt{u}$ is nonincreasing on
$(0,\infty)$. Thus the full RCT-plus-OBS modulus, including its tail
terms, has the sub-root scaling used below.
Thus $u_{j,n_f,N_f}^\star$ is an excess-risk level, not an
$L_2$-distance radius.

\begin{theorem}[Localized two-source excess risk]
\label{thm:cate-general-localized}
Assume RCT consistency and conditional randomization, a known strict-overlap
propensity $e$, honest fixed nuisances, and either a common covariate
marginal with $r_0=1$ or exact transport with
$r_0=dP_R^X/dP_O^X$. Assume the displayed losses are finite, the penalized
objective is bounded below, and a measurable global
$\epsilon_{\mathrm{opt}}$-approximate minimizer exists. Assume the RCT and
OBS tuning samples are independent; the function class and
localized loss-increment classes are measurable or separable; their envelopes
are finite. For every fixed $\delta$, assume the full modulus
$u\mapsto\Phi_{j,n_f,N_f}(u;\delta)$ is nondecreasing and
$u\mapsto\Phi_{j,n_f,N_f}(u;\delta)/\sqrt{u}$ is nonincreasing for
$u>0$. Assume also that, with probability at least
$1-\delta$ for each source,
\begin{align*}
\sup_{\zeta\in\mathcal Z_j(u)}\{-\mathbb Z_{R,j}(\zeta)\}
&\leq\psi_{R,j}(u;n_f)+\operatorname{Tail}_{R,j}(u,\delta),\\
\sup_{\zeta\in\mathcal Z_j(u)}\{-\mathbb Z_{O,j}(\zeta)\}
&\leq\psi_{O,j}(u;N_f)+\operatorname{Tail}_{O,j}(u,\delta)
\end{align*}
simultaneously on the peeling levels used in the proof.  Then, with
probability at least $1-2\delta$,
\begin{align}
\mathcal R(\widehat\zeta_j)-\mathcal R(\zeta_j^\circ)
&\leq C_0\left\{
u_{j,n_f,N_f}^\star(\delta)
+\operatorname{Pen}(\zeta_j^\circ)
+\epsilon_{\mathrm{opt}}\right\},
\label{eq:cate-general-localized}
\end{align}
where $C_0$ depends only on the stated sub-root and fixed-point constants.
The theorem does not assert empirical convexity; the global approximate-ERM
condition remains substantive when the source correction is signed.
\end{theorem}

An estimated density ratio changes both the population loss and the empirical
process.  Condition on the ratio-training sample and let
$u_{j,n_f,N_f}^\star(\widehat r;\delta)$ denote the preceding fixed point
for the $\widehat r$-weighted empirical loss.  Define the fitting drift
\begin{align*}
\Delta_j^{\mathrm{tune}}(u;\widehat r)
&\triangleq
\sup_{\zeta\in\mathcal Z_j(u)}
\left|\{L_j^{\widehat r}-L_j^{r_0}\}(\zeta)
-\{L_j^{\widehat r}-L_j^{r_0}\}(\zeta_j^\circ)\right|.
\end{align*}

\begin{proposition}[General estimated-ratio fitting drift]
\label{prop:cate-general-estimated-r}
Assume RCT consistency and conditional randomization, a known strict-overlap
propensity, honest nuisance fitting, independent RCT and OBS tuning samples,
and an independent ratio-training sample. Suppose either
$P_R^X=P_O^X$ with $r_0=1$, or $P_R^X\ll P_O^X$ with
$r_0=dP_R^X/dP_O^X$. Condition on the ratio-training sample. Assume the
actual $\widehat r$-weighted penalized objective is
bounded below and has a measurable global approximate minimizer, and that its
measurable localized classes have finite envelopes and sourcewise
high-probability modulus bounds. For every fixed $\delta$, assume their full
RCT-plus-OBS modulus $\Phi(u;\delta)$ is nondecreasing and
$\Phi(u;\delta)/\sqrt u$ is nonincreasing for $u>0$, including the stated
tail terms. Form the fixed point from the actual $\widehat r$-weighted loss.
Assume the localization closes at $\bar u$ in the explicit sense that
\begin{align*}
C_1\{u_{j,n_f,N_f}^\star(\widehat r;\delta)
+\Delta_j^{\mathrm{tune}}(\bar u;\widehat r)
+\operatorname{Pen}(\zeta_j^\circ)+\epsilon_{\mathrm{opt}}\}
&\leq\bar u.
\end{align*}
Then, with probability at least $1-2\delta$ conditional on the
ratio-training sample,
\begin{align*}
\mathcal R(\widehat\zeta_j^{\widehat r})-
\mathcal R(\zeta_j^\circ)
&\leq C_1\left\{
u_{j,n_f,N_f}^\star(\widehat r;\delta)
+\Delta_j^{\mathrm{tune}}(\bar u;\widehat r)
+\operatorname{Pen}(\zeta_j^\circ)
+\epsilon_{\mathrm{opt}}\right\}.
\end{align*}
If $\|\widehat r-r_0\|_{L_2(P_O)}\leq a_{n_r,N_r}$ and
\begin{align*}
\sup_{\zeta\in\mathcal Z_j(u)}
\|G_\zeta-G_{\zeta_j^\circ}\|_{L_2(P_O)}
&\leq B_{G,j}(u),
\end{align*}
then
$\Delta_j^{\mathrm{tune}}(u;\widehat r)
\leq|\omega|a_{n_r,N_r}B_{G,j}(u)$.
This is a fitting-drift result, not density-ratio orthogonality.
\end{proposition}

The preceding results control each fixed candidate.  We now select the fusion
coefficients while keeping the final validation observations separate.


% ===== Former Theorem (Risk of the honestly selected general learner), Section 4.4 =====
\begin{theorem}[Risk of the honestly selected general learner]
\label{thm:cate-general-selected-risk}
Let $\widehat\tau_1,\ldots,\widehat\tau_M$ be $M<\infty$ fitted
candidates fixed before final validation and assume
RCT consistency and conditional randomization, known strict-overlap trial
propensity, honest nuisance fitting, independent RCT and OBS tuning samples,
and final RCT and OBS evaluation samples independent of all candidate fitting.
Suppose either the covariate marginals agree and $r_0=1$, or exact transport
holds with $r_0=dP_R^X/dP_O^X$. Allow independently trained ratios in fitting
and validation and condition on them. For candidate $k$, let
$\zeta_k^\circ$ be its best-in-class target, let
$u_k^\star\triangleq
u_{k,n_f,N_f}^\star(\widehat r_k;\delta_f/M)$, and define
\begin{align*}
B_k
&\triangleq\mathcal R(\zeta_k^\circ)
+C_k\{u_k^\star+\Delta_k^{\mathrm{tune}}
+\operatorname{Pen}_k(\zeta_k^\circ)+\epsilon_{\mathrm{opt},k}\}.
\end{align*}
Assume directly that
$\mathcal R(\widehat\tau_k)\leq B_k$ for every
$k=1,\ldots,M$ simultaneously with probability at least
$1-\delta_f$.
Let $\widehat k$ be the deterministically tie-broken index selected by the
common-reference criterion in Equation~\eqref{eq:common-reference-validation}.
Let $\widehat\varepsilon_{\mathrm{eval}}(\delta_v)$ and
$\Delta_r^{\mathrm{eval}}\triangleq\Delta_r$ be the validation concentration
radius and validation ratio drift in
Equation~\eqref{eq:validation-drift-bounds}, respectively. Assume the
validation-selection event
\begin{align*}
\mathcal R(\widehat\tau_{\widehat k})
&\leq \min_{1\leq k\leq M}\mathcal R(\widehat\tau_k)
+2\widehat\varepsilon_{\mathrm{eval}}(\delta_v)
+2\Delta_r^{\mathrm{eval}}
\end{align*}
holds with probability at least $1-\delta_v$. Then, with probability at least
$1-\delta_f-\delta_v$,
\begin{align}
\mathcal R(\widehat\tau_{\widehat k})
&\leq\min_{1\leq k\leq M}B_k
+2\widehat\varepsilon_{\mathrm{eval}}(\delta_v)
+2\Delta_r^{\mathrm{eval}}.
\label{eq:cate-general-selected-risk}
\end{align}
Here $\Delta_k^{\mathrm{tune}}$ changes a fitted candidate, whereas
$\Delta_r^{\mathrm{eval}}$ changes how fixed candidates are ranked.  The
result does not imply coefficient consistency, strict improvement, or
approximation of the continuous coefficient oracle.
\end{theorem}

% ===== Former Proposition (Population-loss flatness), Section 4.4 =====
\begin{proposition}[Population-loss flatness in the fusion scale]
\label{prop:cate-omega-flatness}
Assume RCT consistency and conditional randomization, a known trial
propensity $e$ satisfying strict overlap, and nuisance functions fitted
honestly on data independent of the loss-evaluation observations.  Suppose
either that the RCT and OBS have a common covariate marginal with $r_0=1$,
or that exact covariate transport holds with
$r_0=dP_R^X/dP_O^X$.  Use $r=r_0$, let
$\Omega\subset\mathbb R$ be nonempty and compact, and fix
$\zeta\in\mathcal Z$ and $\lambda\in[0,1]$.  Condition on the fitted
nuisances and suppose the expectations below are finite.  For every
$\omega\in\Omega$, the method-specific corrections satisfy
\begin{align*}
\mathbb E_O\{r_0(X)G_\zeta(X)\}-\mathbb E_R\{G_\zeta(X)\}
&=0,\\
\mathbb E_O\{r_0(X)G_\zeta(X)\}
-\mathbb E_R\{\chi(A,X)G_\zeta(X)\}
&=0.
\end{align*}
Consequently,
\begin{align*}
\mathbb E\{\widehat L_{\mathrm{DRF}}
(\zeta;\lambda,\omega,r_0)\}
&=\mathcal L_{\mathrm{DR}}(\zeta;\lambda),\\
\mathbb E\{\widehat L_{\mathrm{RF}}
(\zeta;\lambda,\omega,r_0)\}
&=\mathcal L_R(\zeta;\lambda),
\end{align*}
and both right-hand sides are constant in $\omega$.  Therefore,
\begin{align*}
\operatorname*{arg\,min}_{\omega\in\Omega}
\mathbb E\{\widehat L_{\mathrm{DRF}}
(\zeta;\lambda,\omega,r_0)\}
&=\Omega,\\
\operatorname*{arg\,min}_{\omega\in\Omega}
\mathbb E\{\widehat L_{\mathrm{RF}}
(\zeta;\lambda,\omega,r_0)\}
&=\Omega.
\end{align*}
Thus population loss at a fixed function and fixed $\lambda$ cannot select
a unique fusion scale.
\end{proposition}

% ===== Former appendix proofs (general two-source learning) =====
\subsection{Proofs for general two-source learning}

For Theorem~\ref{thm:cate-general-erm}, exact transport and the target
identities in Section~\ref{sec:cate-loss} imply
\begin{align*}
L_j(\zeta)-L_j(\zeta_j^\circ)
&=\mathcal R(\zeta)-\mathcal R(\zeta_j^\circ).
\end{align*}
The global approximate-ERM inequality gives
\begin{align*}
\widehat L_j(\widehat\zeta_j)+\operatorname{Pen}(\widehat\zeta_j)
&\leq
\widehat L_j(\zeta_j^\circ)+\operatorname{Pen}(\zeta_j^\circ)
+\epsilon_{\mathrm{opt}}.
\end{align*}
Add and subtract the two population losses.  By the definition of the
centered comparator process, the resulting stochastic term is exactly
$-\mathbb Z_j(\widehat\zeta_j)$.  This proves
Equation~\eqref{eq:cate-general-erm-exact}, including its sign and without an
extra factor two.

For Theorem~\ref{thm:cate-general-localized}, work on the intersection of the
two sourcewise modulus events.  Independence and a union bound give probability
at least $1-2\delta$.  The exact ERM reduction shows that any approximate
minimizer with excess risk in a peeling shell
$2^m u<\mathcal R(\zeta)-\mathcal R(\zeta_j^\circ)\leq2^{m+1}u$
must satisfy an upper bound by the sum of the two source moduli, the penalty at
the comparator, and $\epsilon_{\mathrm{opt}}$.  Sub-root regularity and the
definition of $u_{j,n_f,N_f}^\star$ make that upper bound strictly smaller
than the shell's lower endpoint once $u$ is a sufficiently large universal
multiple of
$u_{j,n_f,N_f}^\star+\operatorname{Pen}(\zeta_j^\circ)
+\epsilon_{\mathrm{opt}}$. For each fixed $\delta$, the full-modulus
sub-root condition gives
$\Phi_{j,n_f,N_f}(tu;\delta)\leq
\sqrt{t}\,\Phi_{j,n_f,N_f}(u;\delta)$ for every $t\geq1$.
Thus the stochastic upper bound grows only as the square root of the peeling
level, whereas the shell's excess-risk lower endpoint grows linearly. Summing the
geometric tail allocation over shells proves
Equation~\eqref{eq:cate-general-localized}.  Separability permits the displayed
suprema to be taken over countable dense subclasses; otherwise the same
argument is understood in outer probability.

For Proposition~\ref{prop:cate-general-estimated-r}, condition on the
ratio-training sample.  The preceding argument now uses the actual
$\widehat r$-weighted source process, hence its fixed point is
$u_{j,n_f,N_f}^\star(\widehat r;\delta)$.  Add and subtract the exact-ratio
population loss at $\zeta$ and $\zeta_j^\circ$.  Their increment is bounded
by $\Delta_j^{\mathrm{tune}}$, after which the same localization argument
closes at $\bar u$.  If the two stated $L_2(P_O)$ bounds hold,
Cauchy--Schwarz gives
\begin{align*}
\Delta_j^{\mathrm{tune}}(u;\widehat r)
&\leq |\omega|\|\widehat r-r_0\|_{L_2(P_O)}
\sup_{\zeta\in\mathcal Z_j(u)}
\|G_\zeta-G_{\zeta_j^\circ}\|_{L_2(P_O)},
\end{align*}
which proves the final claim.


% ===== Former proof of the selected-learner theorem =====
For Theorem~\ref{thm:cate-general-selected-risk}, assign fitting failure
probability $\delta_f/M$ to each of the $M$ fixed candidates.  A union
bound makes all candidate-specific fitting inequalities hold simultaneously
with probability at least $1-\delta_f$.  On that event, the ideal risk of
candidate $k$ is at most $B_k$.  Independently, the validation event in
Corollary~\ref{cor:cate-estimated-r-validation} holds with probability at
least $1-\delta_v$ and compares the selected fitted candidate with the
best fitted candidate, adding
$2\widehat\varepsilon_{\mathrm{eval}}(\delta_v)
+2\Delta_r^{\mathrm{eval}}$.  Intersecting the fitting and validation events
and then minimizing the simultaneous upper bounds over $k$ proves
Equation~\eqref{eq:cate-general-selected-risk}.  No independence among the
candidate fits is required; only their joint independence from final
validation is used.

% ===== Former proof of the flatness proposition =====
For Proposition~\ref{prop:cate-omega-flatness}, exact transport gives
\begin{align*}
\mathbb E_O(r_0G_\zeta)-\mathbb E_R(G_\zeta)=0.
\end{align*}
Moreover, conditional randomization and strict overlap give
$\mathbb E_R(\chi\mid X)=1$, so
\begin{align*}
\mathbb E_R(\chi G_\zeta)
&=\mathbb E_R[\mathbb E_R(\chi\mid X)G_\zeta]
=\mathbb E_R(G_\zeta),\\
\mathbb E_O(r_0G_\zeta)-\mathbb E_R(\chi G_\zeta)
&=0.
\end{align*}
Substitution into Equations~\eqref{eq:dr-fusion-loss} and
\eqref{eq:rf-fusion-loss} leaves
$\mathcal L_{\mathrm{DR}}(\zeta;\lambda)$ and
$\mathcal L_R(\zeta;\lambda)$, respectively, for every
$\omega\in\Omega$.  Each objective is therefore constant on the nonempty
set $\Omega$, and its minimizer set is all of $\Omega$.
